\section{Projects}

\newcommand{\smallsep}[0]{\vspace{0.2em}}

    % \runsubsection{Twitter-Emotion~~|} \descript{Tweet emotion classifier in Pytorch
    % \href{https://github.com/zoravur/twitter-emotion/blob/master/Twitter\%20Emotion\%20Classification.ipynb}{\faGithub}}
    % \smallsep
    % \begin{tightemize}
    %     \item \textbf{Trained a CNN} to classify the emotional valence of a given tweet and achieved an F1 score of 81\%.
    %     \item Used GloVe \textbf{word vectors} trained on a Twitter dataset of over 2 billion tweets, resulting in a \textasciitilde20\% increase in accuracy over a bag-of-words approach.
    % \end{tightemize}
    % \sectionsep
    % \runsubsection{Gravity.js~~|} \descript{A highly accurate HTML5 gravity simulation
    % \href{https://github.com/zoravur/gravity}{\faGithub}
    % \href{https://zoravur.github.io/gravity/dist}{\faLink \enskip \md https://zoravur.github.io/gravity/dist}}
    % \smallsep
    % \begin{tightemize}
    %     \item Created an interactive newtonian gravity simulation in TypeScript, designed for exploring orbital mechanics.
    %     % \item Implemented the semi-implicit euler method, in order to ensure stability of a long period of time.
    %     \item Improved performance benchmark from 20FPS to 57FPS by using multithreading for physics calculations.
    %     % \textbf{numerical integration methods}, and an intuitive UI for controlling simulation parameters.
    %     % \item Used \textbf{multithreading} through Web Workers and HTML5 Canvas to achieve optimal performance.
    %     % Created using NPM, webpack,
    % \end{tightemize}
    % \sectionsep 

    \runsubsection{Spock~~|} \descript{A spoken programming language
    \href{https://github.com/zoravur/spock}{\faGithub}
    \href{https://zoravur.com/spock-a-spoken}{https://zoravur.com/spock-a-spoken}}
    \smallsep
    \begin{tightemize}
        \item Created a tree-walk interpreter for a spoken programming language consisting of english words and no punctuation.
        % \item Implemented the semi-implicit euler method, in order to ensure stability of a long period of time.
        \item Used OpenAI's STT model Whisper to dictate a complete solution to Project Euler problem 1.
        % \textbf{numerical integration methods}, and an intuitive UI for controlling simulation parameters.
        % \item Used \textbf{multithreading} through Web Workers and HTML5 Canvas to achieve optimal performance.
        % Created using NPM, webpack,
    \end{tightemize}
    \sectionsep 

    \runsubsection{Taan Generator~~|} \descript{A spoken programming language
    \href{https://github.com/zoravur/taangenerator}{\faGithub https://github.com/zoravur/taangenerator}}
    \smallsep
    \begin{tightemize}
        \item Trained a transformer to generate taans (runs of musical notes) in the Indian classical raga (mode) \href{https://en.wikipedia.org/wiki/Yaman_(raga)}{Yaman}. \href{https://https://www.youtube.com/watch?v=FXV69hGpr4Q}{\color{myyellow} \selectfont{(demo)}}
        % \item Implemented the semi-implicit euler method, in order to ensure stability of a long period of time.
        % \textbf{numerical integration methods}, and an intuitive UI for controlling simulation parameters.
        % \item Used \textbf{multithreading} through Web Workers and HTML5 Canvas to achieve optimal performance.
        % Created using NPM, webpack,
    \end{tightemize}
    \sectionsep     
    % \runsubsection{Stream-transduce~~|} \descript{Library for functional programming with Node streams
    % \href{https://github.com/zoravur/stream-transduce}{\faGithub}
    % \href{https://www.npmjs.com/package/stream-transduce}{\faNpm}}
    % \smallsep
    % \begin{tightemize}
    %     \item Extended functionality of the existing functional programming library Ramda by creating a functional wrapper for Node.js stream objects.
    %     \item Published to NPM and ensured usability through documentation (with JSDoc) and semantic versioning.
    % \end{tightemize}
    
    % \runsubsection{VM~~|}
    % \descript{Vim text editor clone in C++}
    % \smallsep
    % \begin{tightemize}
    %     % \item Developed a modal text editor with 55 Vim commands, syntax highlighting, and recursive macros.
    %     % % \item Created a fully featured clone of the text editor Vim in C++, with support for normal and insert mode, common commands, and macros.
    %     % % \item Created abstractions for the a wrapper C-lev input library ncurses
    %     % \item Used Flux, reactive programming, and various Object Oriented Design patterns.
    %     % \item Received a mark of \textbf{96\%} (class average 70\%).
    %     \item Created a modal text editor supporting 55 vi commands, normal and insert mode, and macros.
    %     \item Received a mark of 96\% as the final project for the Object Oriented Design course.
    % \end{tightemize}
    % \sectionsep
    
    % \runsubsection{Visually Study Yo Code~~|} \descript{Data structure visualization debugger extension \href{https://github.com/cherrykit/VSYCode}{\faGithub}}
    % \smallsep
    % \begin{tightemize}
    %     \item Created a VSCode extension to visualize the current state of variables during a debug session.
    %     \item Implemented custom visualization functions for binary trees, heaps, linked lists, and arrays.
    %     \item Implemented real time updating of graphics as users stepped through the debug session.
    % \end{tightemize}
    % \sectionsep 
    
    % \runsubsection{Serv~~|} \descript{A collaborative reporting app for the Toronto Police}
    % \smallsep
    % \begin{tightemize}
    %     \item Built a Firebase application that facilitated the sharing of video evidence of crimes.
    %     \item Implemented face detection in Python using \textbf{OpenCV} to detect potential suspects in security footage.
    %     % \item \textbf{Tech Stack:} Front End: HTML, CSS, JS; Backend: Python, Node.js, Firebase; Used Google Maps API, OpenCV, and other Web APIs.
    %     % \item Used web APIs and front-end libraries to create an aesthetic and functional UI.
    %     \item Submitted at the hackathon HackTPS.
    % \end{tightemize}
    % \sectionsep


% }{


    % \runsubsection{SpecZo}
    % \descript{News bias reporter}
    % %\location{Hack the North, Waterloo}
    % \begin{tightemize}
    %     \item Built a chrome extension designed to report bias of different news sources for the Hackathon Hack the North.
    %     \item Accessed chrome data, etc.
    %     \item something else impressive haha
    % \end{tightemize}

% }
    



% \runsubsection{Votebot}
% \descript{}
% \location{2nd Prize, Hackademics VN 2015}
% \item Deployed nodejs + puppeteer bot designed to vote consistently
% A pioneering sex-ed website that promotes sexuality and gender awareness among Vietnamese teenagers. Built with MEAN stack.
% \sectionsep

% \begin{itemize}
% \item Built a full-stack web application that facilitated the sharing of evidence related to crime (designed for the Toronto Police).
% \item Used the computer vision library OpenCV to recognize faces in large video datasets.
% \item \textbf{Tech Stack:} Front End: HTML, CSS, JS; Backend: Python, Node.js, Firebase; Used Google Maps API, OpenCV, and other Web APIs.
% % \item Used web APIs and front-end libraries to create an aesthetic and functional UI.
% %\item Submitted at the hackathon HackTPS. 



    % \runsubsection{Serv} \descript{A collaborative reporting app}
    % \begin{tightemize}
    %     \item Built a web application that facilitated the sharing of evidence related to crime (designed for the Toronto Police).
    %     \item Used the computer vision library OpenCV to recognize faces in large video datasets.
    %     % \item \textbf{Tech Stack:} Front End: HTML, CSS, JS; Backend: Python, Node.js, Firebase; Used Google Maps API, OpenCV, and other Web APIs.
    %     % \item Used web APIs and front-end libraries to create an aesthetic and functional UI.
    %     \item Submitted at the hackathon HackTPS.
    % \end{tightemize}
    % \sectionsep

%}{